\section{Introduction}
\label{sec:introduction}

Pengenalan otomatis atribut ras dan gender wajah memegang peranan krusial pada tata kelola identitas digital, kontrol akses biometrik, forensik digital, dan personalisasi antarmuka cerdas~\cite{ref1}, \cite{ref2}. Berbagai studi independen dan evaluasi National Institute of Standards and Technology (NIST) menunjukkan bahwa model visi komputer komersial memperlihatkan disparitas performa yang merugikan kelompok minoritas tertentu, khususnya wanita berkulit gelap~\cite{ref3}, \cite{ref4}. Kesenjangan akurasi ini bersumber dari karakteristik distribusi data latih yang tidak seimbang serta representasi fitur visual yang kurang tangguh terhadap variasi fenotipik dan pencahayaan lingkungan~\cite{ref5}. Keterbatasan sistem konvensional dalam menangani variasi visual tersebut memperlihatkan perlunya pendekatan representasi baru yang memodelkan interaksi ras dan gender secara terpadu.

Praktik umum dalam literatur visi komputer memisahkan klasifikasi ras dan gender sebagai dua model komputasi independen yang berdiri sendiri~\cite{ref6}. Pendekatan terisolasi tersebut gagal menangkap dinamika interaksi dan kesalahan pada persilangan subkelompok demografis tertentu, seperti Black Females atau Asian Females~\cite{ref7}, \cite{ref8}. Akibatnya, penurunan performa klasifikasi pada irisan demografis spesifik tersebut sering kali luput dari audit metrik akurasi agregat konvensional~\cite{ref9}. Formulasi interseksional terpadu yang memetakan enam kelas gabungan (3 Ras × 2 Gender) pada kondisi evaluasi seimbang memungkinkan audit performa yang adil dan transparan. Pemodelan klasifikasi interseksional ini menuntut representasi visual kaya yang mampu memisahkan batas keputusan antarsubkelompok yang saling tumpang tindih secara fenotipik~\cite{ref10}.

Metodologi pengenalan atribut wajah telah berevolusi dari fitur buatan tangan klasik menuju arsitektur Convolutional Neural Network (CNN) dan Vision Transformer (ViT). Pada era CNN, penelitian terdahulu mengeksplorasi ekstraksi ciri demografis melalui pemanfaatan wilayah wajah tengah~\cite{ref11} serta kerangka kerja multi-tugas multi-skala untuk estimasi usia dan gender~\cite{ref12}. Di samping itu, model CNN mendalam turut dikembangkan untuk meningkatkan akurasi pengenalan etnisitas berbasis citra wajah~\cite{ref13}. Arsitektur ViT kemudian hadir membawa keunggulan pemodelan hubungan spasial global antarpatch citra melalui mekanisme Multi-Head Self-Attention (MHSA)~\cite{ref14}. Pendekatan representasi generatif sintetis juga dieksplorasi untuk memitigasi disparitas demografis lintas kelompok gender dan ras~\cite{ref15}. Perkembangan mutakhir mencakup integrasi modul atensi paralel~\cite{ref16} serta pemanfaatan ViT untuk klasifikasi gender wajah~\cite{ref17}. Selain itu, eksplorasi arsitektur transformer multi-aksis diusulkan untuk memperkaya variasi pemodelan representasi visual hierarkis~\cite{ref18}. Penggabungan representasi ViT lintas domain mulai dirintis melalui integrasi fitur ganda wajah dan emosi~\cite{ref19}, serta fusi representasi multi-domain untuk pengenalan atribut demografis yang adil~\cite{ref20}. Terlepas dari kemajuan pesat model visi transformer tersebut, telaah literatur mengidentifikasi sejumlah tantangan metodologis yang masih memerlukan kajian empiris terpadu.

Keterbatasan representasi domain tunggal menjadi kesenjangan metodologis utama karena biometrik wajah statis rentan terhadap distorsi ekspresi afektif dinamis~\cite{ref11}, \cite{ref17}. Selain itu, sistem pengenalan wajah sensitif terhadap variasi penuaan biologis jika fitur morfologi usia diabaikan~\cite{ref12}, \cite{ref18}. Kesenjangan kritis lainnya adalah ancaman kebocoran identitas (identity leakage), di mana evaluasi konvensional berbasis partisi acak citra (image-level random split) membagi foto dari subjek yang sama ke data latih dan uji, sehingga memicu memorisasi identitas (shortcut learning). Terlebih lagi, studi terdahulu umumnya minim pembuktian signifikansi inferensial formal~\cite{ref13}, \cite{ref19}. Profil efisiensi komputasi praktis pada pengklasifikasi hilir juga belum dieksplorasi secara mendalam~\cite{ref20}. Untuk menjembatani kesenjangan tersebut, penelitian ini mengusulkan fusi tri-domain ViT di bawah protokol evaluasi ketat.

Untuk menjawab kesenjangan tersebut, penelitian ini mengusulkan fusi Tri-Domain ViT yang memadukan geometri biometrik statis (ViT-Face), ekspresi afektif dinamis (ViT-Emotion), dan morfologi penuaan biologis (ViT-Age) melalui operasi konkatenasi langsung (Face $\oplus$ Emotion $\oplus$ Age: 2.304 dimensi)~\cite{ref21}, \cite{ref22}. Ekstraksi representasi laten visual dilakukan secara offline dari backbone ViT pra-latih yang dibekukan (frozen) untuk menjaga efisiensi komputasi tanpa fine-tuning ulang~\cite{ref23}. Penelitian ini mengadopsi arsitektur evaluasi Dual-Benchmark, yang menetapkan Strict Subject-Disjoint Partitioning (600 subjek unik, 480 latih dan 120 uji held-out, 0\% kebocoran) sebagai tolok ukur utama generalisasi, sekaligus mempertahankan evaluasi subject-overlap historis sebagai studi ablasi empiris untuk menguantifikasi efek memorisasi identitas. Optimasi hyperparameter komprehensif dilakukan pada empat famili pengklasifikasi klasik, yaitu Random Forest (RF), Gaussian Naive Bayes (GNB), Logistic Regression (LR), dan Support Vector Machine (SVM), melalui validasi silang Cross-Validation (CV) 5-Fold StratifiedGroupKFold. Pipeline modular penskalaan dan Principal Component Analysis (PCA) difitting secara terisolasi pada lipatan latih untuk mencegah kebocoran informasi dengan total 31.920 penyesuaian model~\cite{ref24}.

The primary scientific contributions of this research are summarized into four core aspects as follows (Kontribusi ilmiah utama dari penelitian ini dirangkum ke dalam empat aspek esensial sebagai berikut):
\begin{itemize}
\item \textbf{A Systematic Multi-Domain ViT Latent Fusion \& Ablation Framework:} Formulating a seven-scheme feature ablation framework that integrates frozen task-specific Vision Transformer backbones (craniofacial biometrics, affective expressions, and biological age) with leakage-free fold-level preprocessing and hyperparameter optimization across four classifier families for intersectional demographic classification. (Formulasi kerangka kerja ablasi tujuh skema fitur yang mengintegrasikan backbone Vision Transformer pra-latih yang dibekukan dengan prapemrosesan tingkat-fold bebas kebocoran serta optimasi hyperparameter lintas empat famili pengklasifikasi untuk klasifikasi demografis interseksional).
\item \textbf{Strict Subject-Disjoint Benchmark \& Data Leakage Ablation:} Establishing a zero-leakage benchmark on the DemogPairs dataset and providing a comparative ablation study to quantify the facial identity memorization effect inherent in conventional random splits. (Penetapan tolok ukur evaluasi bebas kebocoran identitas pada dataset DemogPairs dan studi komparatif ablasi kebocoran data untuk menguantifikasi efek memorisasi identitas wajah).
\item \textbf{Inferential Hypothesis Testing \& Classifier-Dependent Saturation Characterization:} Conducting rigorous paired inferential testing to validate multi-domain gains over single-domain baselines, showing no statistically significant predictive difference between top linear and non-linear classifiers, and characterizing linear age-feature complementarity versus non-linear asymptotic representation saturation across classifier families. (Pengujian hipotesis inferensial berpasangan yang memvalidasi peningkatan signifikan representasi multi-domain atas baseline domain tunggal, menunjukkan tidak adanya perbedaan performa prediktif yang signifikan secara statistik antara pengklasifikasi linier dan non-linier terbaik, serta mengkarakterisasi komplementaritas fitur usia linier versus kejenuhan representasi asimtotik non-linier lintas famili pengklasifikasi).
\item \textbf{Comprehensive Computational Efficiency \& Downstream Classifier Profiling:} Delivering a comprehensive computational profiling of ViT extraction time, training cost, inference latency, and memory footprint, establishing an ultra-lightweight downstream classification head suitable for practical deployment. (Penyajian profil efisiensi komputasi komprehensif mencakup waktu ekstraksi ViT, biaya pelatihan, latensi inferensi, dan jejak memori yang merealisasikan kepala pengklasifikasi hilir berbobot sangat ringan untuk implementasi praktis).
\end{itemize}

Sistematika penulisan artikel ini disusun sebagai berikut: Section~\ref{sec:related_works} mengulas sintesis literatur terkait pengenalan demografis wajah dan arsitektur transformer. Section~\ref{sec:materials_and_methods} menjabarkan spesifikasi dataset, metodologi ekstraksi fitur multi-domain ViT, serta pipeline optimasi pengklasifikasi. Section~\ref{sec:results_and_discussion} menyajikan hasil evaluasi eksperimental, studi ablasi kebocoran data, pengujian inferensial, profil efisiensi komputasi, dan analisis keadilan interseksional. Terakhir, Section~\ref{sec:conclusion} merangkum kesimpulan utama, batasan metodologis, serta rekomendasi arah riset lanjutan.
